\documentclass[10pt]{article}
\usepackage[margin=20mm]{geometry}
\usepackage[T1]{fontenc}
\usepackage{lmodern}
\usepackage{hyperref}
\begin{document}
\section*{Positive shared-support distinct-other: word composition}
\textbf{Component result.} A literal UInt64 value and conditional-dispatch composition connects the reviewed mixed prelude, degree-one doubling and generic mixed-addition implementations. The theorem derives nonzero original ordinate, a successful Some return, canonical output, nonzero output weights and normalized Mumford validity from original input validity. No output-validity assumption is introduced.

\begin{center}\fbox{Canonical valid original inputs}\par $\downarrow$\par
\fbox{Outer $K=0$, cancellation false, distinct other support}\par $\downarrow$\par
\fbox{Reviewed field branch derives positive equality and $y\ne0$}\par $\downarrow$\par
\fbox{Literal word doubling $\rightarrow$ generic word dispatch}\par $\downarrow$\par
\fbox{Some canonical, nonzero-weight, normalized-valid output}\end{center}

\textbf{Premises.} Prime $5\le p<2^{31}$; canonical ten input coordinates and four curve coefficients; four nonzero input weights; normalized degree-one and degree-two divisibility on $X^5+a_3X^3+a_2X^2+a_1X+a_0$. The three branch premises are literal word tests: outer $K=0$, inverse/cancellation comparison false, and the cross-scaled original coordinate unequal to the extracted other coordinate. Positive equality is derived.

\textbf{Source mapping.} Retained Rust add12\_weighted\_complete lines828--878, doubling477--510 and generic232--292. The xz2 multiplication tree is literal; existing WordMixed, WordDoubling and WordGeneric provide payload decode and canonicality. The conditional word wrapper returns None outside the scoped positive-distinct route as a restriction marker; it does not represent the other production branches. Nonzero original word y justifies the degree-two doubling extraction. The value mapping uses Field.n=1 and tables=None; counter state and compiler execution are outside this theorem.

\textbf{Checks.} Main module attempt7 and separate axiom readback use the existing pinned Lean environment with a120-second command bound. Readback lists only propext, Classical.choice and Quot.sound. The full nonunit p=7 fixture uses a=(0,1,2,2), b=(5,0,0,3,3,2), coefficients(1,6,0,0), with explicit polynomial divisibility witnesses for both original inputs. Additional controls cover cancellation precedence, repeated-other exclusion and noncanonical input rejection. Native receipt status is recorded in status.json; independent source review follows sealing.

\textbf{Retained failures.} Early attempts exposed a let-bound split issue, heartbeat exhaustion and two120-second caps caused by elaborating an underspecified normalized-validity conversion. An explicit abstract transport lemma avoids re-expanding nested arithmetic. A missing prime instance in that helper was corrected. Fixture attempts retain finite-field simplification and decidability errors; the assumptions and fixture were preserved.

Repeated-support tripling, full mixed dispatch, divisor-class addition, whole-walk closure, counters and compiler behavior are separate obligations. No scientific trials or timing measurements were run; cost ratio is null.
\end{document}
